\section{Falsifiable Predictions and Empirical Research Pathway}
\label{sec:hypotheses}

Although \RST{} is framed around a satirical premise, its structural mechanisms make
concrete empirical claims whenever real-world relational portfolios behave analogously
to the model. Because Lean proofs and synthetic simulations can only verify mathematical
consistency, all behavioral propositions remain empirically unestablished until tested
against observational data. Below we formulate four falsifiable empirical hypotheses, each
following an explicit eight-part specification.

\subsection{Hypothesis H1: Option-Supply Mediation of Concurrent Portfolios}
\label{subsec:hyp-h1}

\noindent\layerempiric{}
\begin{description}[leftmargin=1.4em,nosep]
  \item[Hypothesis statement:] Equilibrium prevalence of sustained multi-partner portfolios
        ($n_{\mathrm{eff}} \ge 2$) increases monotonically with perceived attractiveness $A$,
        and this association is predominantly mediated by viable incoming option count $n(A)$.
  \item[Construct operationalization:] Attractiveness $A_i$ via out-of-sample panel ratings
        (standardized composite); $N_i$ via $14$-day incoming unique dating interest
        (\texttt{OPT\_INCOMING}); superposition $Y_{it} = \mathbf{1}\{\texttt{W\_ACTIVE\_N}_{it} \ge 2\}$
        for $\ge 3$ consecutive weeks (\cref{app:instrument}).
  \item[Null hypothesis:] Neither $\beta_1$ (unadjusted $A_i$ slope) nor $\beta_2$ (mediated
        $\log(1 + N_i)$ slope) in~\eqref{eq:instrument-regression} differs from zero.
  \item[Minimal study design:] Longitudinal $8$-week diary cohort ($N_{\mathrm{subj}} \ge 400$
        single adults) with weekly passive/diary interaction logs avoiding explicit ranking prompts.
  \item[Exclusion and bias risks:] Social desirability under-reporting of multiple chats;
        attrition skewed toward newly committed participants exiting early.
  \item[Confounders:] Sociosexual orientation, attachment avoidance, geographic density,
        and endogenous selective filtering (high-$A$ selectivity offsetting higher $N_i$).
  \item[Confirmation criteria:] Significant $\beta_2 > 0$ ($p < 0.01$) alongside sharp
        attenuation ($>50\%$) of direct coefficient $\beta_1$ once $\log(1 + N_i)$ is entered.
  \item[Falsification conditions:] Zero or negative partial association ($\beta_2 \le 0$)
        or persistent strong $\beta_1 > 0$ wholly unmediated by incoming option volume $N_i$.
\end{description}

\subsection{Hypothesis H2: Query-Order Asymmetry Under Direct Status Inquiries}
\label{subsec:hyp-h2}

\noindent\layerempiric{}
\begin{description}[leftmargin=1.4em,nosep]
  \item[Hypothesis statement:] Asking ``where do I rank?'' prior to ``what functional role
        do I fill?'' yields systematically different classification stability and commitment
        shifts than asking role before rank (\cref{subsec:poke-sequential}).
  \item[Construct operationalization:] Randomized vignette or behavioral audit assignment to
        order $R \to L$ vs.\ $L \to R$, measuring post-prompt categorical agreement and
        response latency.
  \item[Null hypothesis:] Joint conditional answer distributions match within sampling noise
        ($\Prob(r,\ell \mid R\to L) = \Prob(r,\ell \mid L\to R)$).
  \item[Minimal study design:] Randomized two-arm survey experiment ($N_{\mathrm{subj}} \ge 300$)
        among respondents currently managing $\ge 2$ ambiguous romantic chats.
  \item[Exclusion and bias risks:] Demand characteristics and hypothetical vignette detachment.
  \item[Confounders:] Prior conversation salience, baseline ambiguity tolerance, recency
        of last contact.
  \item[Confirmation criteria:] Statistically significant Jensen-Shannon divergence between
        $R\to L$ and $L\to R$ joint contingency tables ($p < 0.01$ via permutation test).
  \item[Falsification conditions:] Symmetric joint response distributions across both query
        orders consistent with commuting classical evaluation.
\end{description}

\subsection{Hypothesis H3: Social-Graph Density Amplifies Simultaneous Dissolution}
\label{subsec:hyp-h3}

\noindent\layerempiric{}
\begin{description}[leftmargin=1.4em,nosep]
  \item[Hypothesis statement:] Holding active partner count $n$ fixed, higher pairwise friend
        edge density $q$ increases the hazard of simultaneous multi-partner exit within $14$ days
        of a conflictual poke (\cref{sec:cascade}).
  \item[Construct operationalization:] Baseline network density $\bar{q}_i$ computed from
        \texttt{NET\_OVERLAP}; simultaneous dissolution coded as $\ge 2$ partners dropping to
        zero weekly contact inside $14$ days following any \texttt{P\_POKE\_EVENT}.
  \item[Null hypothesis:] Edge density coefficient $\beta_3 = 0$ after adjusting for $n$.
  \item[Minimal study design:] Sub-cohort event-history analysis tracking post-poke $14$-day
        retention windows conditional on baseline overlap matrix $G_{\mathrm{soc}}$.
  \item[Exclusion and bias risks:] Hidden digital cross-talk via anonymous social feeds
        unmeasured in self-reported acquaintance matrices.
  \item[Confounders:] Partner impulsivity, shared workplace norms, visibility of public outings.
  \item[Confirmation criteria:] Positive hazard ratio for $\bar{q}_i$ predicting multi-partner
        dissolution ($\ge 2$ exits) over single-partner exit.
  \item[Falsification conditions:] Multi-partner exits occurring independently of measured
        friend overlap ($\beta_3 \approx 0$) or decreasing with shared social embeddedness due
        to reputational coordination.
\end{description}

\subsection{Hypothesis H4: Regime Switching Dissociates Messaging Share from Stated Rank}
\label{subsec:hyp-h4}

\noindent\layerempiric{}
\begin{description}[leftmargin=1.4em,nosep]
  \item[Hypothesis statement:] Day-of-week and stress-period transitions create systematic
        rank reversals in empirical messaging share relative to self-reported priority
        (\cref{sec:dynamics}).
  \item[Construct operationalization:] Passive timestamps grouped into weekday, weekend, and
        self-reported stress blocks vs.\ baseline importance scores; entropy $H(\mathbf{w}(t))$
        per $14$-day rolling bin.
  \item[Null hypothesis:] Conditional partner share $w_i(t \mid r)$ is invariant across regimes $r$.
  \item[Minimal study design:] Longitudinal time-series panel comparing regime-conditioned
        transition frequencies against homogeneous null transition matrices.
  \item[Exclusion and bias risks:] Missing offline interactions distorting message-count weights.
  \item[Confounders:] Partner availability schedules (shift work) rather than focal agent preference.
  \item[Confirmation criteria:] Rejection of regime-homogeneity ($\chi^2$ test across transition
        tables) matching qualitative role-regime specialization.
  \item[Falsification conditions:] Constant transition shares across weekday, weekend, and
        stress windows.
\end{description}
